\documentclass[14pt,a4paper]{extarticle}

\usepackage[T1,T2A]{fontenc}
\usepackage[utf8]{inputenc}
\usepackage[english,russian]{babel}

\usepackage{amsmath,amssymb,amsfonts}
\usepackage{geometry}
\usepackage{setspace}
\usepackage{indentfirst}
\usepackage{longtable}
\usepackage{booktabs}
\usepackage{array}
\usepackage{caption}
\usepackage{hyperref}
\usepackage{enumitem}
\usepackage{titlesec}

\usepackage{tempora} 
\usepackage{newtxmath}

\geometry{left=3cm,right=1.5cm,top=2cm,bottom=2cm}
\setstretch{1.5}
\setlength{\parindent}{1.25cm}
\setlength{\parskip}{0pt}

\hypersetup{
colorlinks=true,
linkcolor=black,
urlcolor=black,
citecolor=black
}

\titleformat{\section}{\bfseries\large}{\thesection.}{0.5em}{}
\titleformat{\subsection}{\bfseries\normalsize}{\thesubsection.}{0.5em}{}

\usepackage{tocloft}
\renewcommand{\cftsecleader}{\cftdotfill{\cftdotsep}} 


\usepackage{graphicx}
\usepackage{tabularx}
\begin{document}
\sloppy
\thispagestyle{empty}
\begin{center}
Национальный исследовательский ядерный университет «МИФИ»
\vspace{3cm}

{\bfseries ОТЧЁТ}
\vspace{1cm}

{\bfseries Название исследования}
\end{center}
\vfill
\begin{flushright}
Выполнил: Зюков А. А.
\end{flushright}
\vfill
\begin{center}
Город, год
\end{center}
\newpage
\tableofcontents
\newpage
\section{Постановка задачи}
Предмет, вопрос, сравниваемые варианты и ожидаемый результат.
\section{Исходные данные и методика}
Источник данных и условия опыта~\cite{manual}.
\begin{equation}
U=\sum_i F_i(\rho_i)+\frac12\sum_{i\ne j}\phi_{ij}(r_{ij}).
\label{eq:energy}
\end{equation}
Обозначения поясняются в тексте после формулы.
\section{Результаты}
В таблице~\ref{tab:settings} приведены настройки.
\begin{table}[htbp]
\centering
\caption{Пример оформления настроек}
\label{tab:settings}
\small
\begin{tabularx}{\textwidth}{@{}Xrr@{}}
\toprule
Вариант & Параметр, эВ & Расстояние, \AA\\
\midrule
Исходный & 1,00 & 2,00\\
Изменённый & 1,25 & 2,00\\
\bottomrule
\end{tabularx}
\end{table}
\section{Выводы}
Ответ на вопрос исследования и область применимости.
\begin{thebibliography}{99}
\addcontentsline{toc}{section}{Список литературы}
\bibitem{manual} Автор. Название. Год. URL: \url{https://example.org}.
\end{thebibliography}
\end{document}
